\subsection{正埃尔米特型。}

定义 1. --- 称 E 上的埃尔米特型 \(\Phi\) 是正的（相应地，负的），如果 \(\Phi(x, x) \geqslant 0\) （相应地，\(\Phi(x, x) \leqslant 0\) ）对一切 \(x \in \mathbf{E}\) 成立。

当 A = K 时，也称 \(Q(x) = \frac{1}{2} \Phi(x, x)\) （与 \(\Phi\) 相伴的二次型）是正的（相应地，负的）。

设 E 在 A 上有限维，且设 \((e_{i})\) ( \(i = 1, \ldots, n\) ) 是 E 的正交基（ \(\S 6, no. 1\) ，定理 1）。要使 E 上的埃尔米特型 \(\Phi\) 是正的，必须且只需 \(\Phi(e_{i}, e_{i}) \geqslant 0\) （对 \(i = 1, \ldots, n\) ）。设 \(\Phi\) 是 E 上的一个非退化正埃尔米特型；由于 K 的每个正元素都是平方，因而具有形式 \(\rho\bar{\rho} (\rho \in A)\) ，故 E 中存在关于 \(\Phi\) 的标准正交基（ \(\S 6, no. 1\) ，定理 1 的推论 1）。

命题 1. --- 假设 E 有限维，且 A = K 或 A = K(i)。若 \(\Phi\) 是 E 上的非退化正埃尔米特型，则 \(\Phi\) 到 \(\bigotimes^{p}\mathbf{E}\) 与 \(\bigwedge^{p}\mathbf{E}\) ( \(p > 0\) ) 上的扩张以及 \(\Phi\) 的逆型都是非退化正埃尔米特型。

这由 E 的标准正交基的存在性及 § 6 第 1 小节命题 2 立即得出。

若把各处的“非退化正”都换成“正”，命题 1 仍成立。

命题 2. --- 设 Φ 是 E 上的一个正埃尔米特型。对 E 中的 x、y，有

\[
\Phi (x, y) \overline {{{{\Phi (x , y)}}}} \leqslant \Phi (x, x) \Phi (y, y).\tag{1}
\]

当向量 \(x\) 与 \(y\) 成比例时，不等式显然是直接的。因此设 \(x\) 与 \(y\) 线性无关。设 A′ 是 A 的（交换）子域 K( \(\Phi(x, y)\) )，F 是 A′ 上由 \(x\) 与 \(y\) 生成的向量平面，\(\Phi_{\mathbb{F}}\) 是 \(\Phi\) 在 F 上的限制；它取值于 A′ 中。由命题 1，\(\Phi_{\mathbb{F}}\) 关于基 \((x, y)\) 的判别式 \(\geqslant 0\) 。而这个判别式就是 \(\Phi(x, x)\Phi(y, y) - \Phi(x, y)\overline{\Phi(x, y)}\) 。证毕。

推论. --- E 的迷向向量所成的集是子空间 \(E^{0}\) （E 关于 \(\Phi\) 的正交子空间）。要使 \(\Phi\) 非退化，必须且只需 \(\Phi(x, x) > 0\) （对一切 \(x \neq 0\) ）。

命题 3. --- 假设 E 有限维且 A 交换（A = K 或 A = K(i)）。设 \(\Phi\) 是 E 上的埃尔米特型，X 是它关于 E 的基 ( \(x_j\) ) ( \(j = 1, \ldots, n\) ) 的矩阵；对 {[}1, n{]} 的每个子集 H，以 \(X_{\mathrm{H,H}}\) 表示从 X 中删去指标 \(j \notin \mathrm{H}\) 的行与列所得的 X 的子式（第 III 章 § 6 第 3 小节）。

\begin{enumerate}
\def\labelenumi{\alph{enumi})}
\item
  若 \(\Phi\) 非退化且正，则 \(X_{\mathrm{H,H}} > 0\) （对 \([1, n]\) 的每个子集 H）。
\item
  反之，若 \(H_{j} = (1, j)\) 且 \(X_{H_{j}, H_{j}} > 0\) （对 \(j = 1, \ldots, n\) ），则 \(\Phi\) 非退化且正。
\end{enumerate}

先设 \(\Phi\) 非退化且正。诸元素 \((x_{j})\) ( \(j\in\mathbb{H}\) ) 构成 E 的一个子空间 F 的基，而子式 \(X_{\mathrm{H,H}}\) 是限制型 \(\Phi_{\mathrm{F}}\) （即 \(\Phi\) 在 F 上的限制）关于这个基的判别式；今因 \(\Phi_{\mathrm{F}}(x,x)>0\) （对 F 中一切 \(x\neq0\) ），故 \(\Phi_{\mathrm{F}}\) 非退化且正（命题 2 的推论）；因而有 \(X_{\mathrm{H,H}}>0\) （命题 1）。为证明 \(b\) )，注意按 § 6 第 1 小节命题 1 的记号，子式 \(X_{\mathrm{Hj,Hj}}\) 等于 D \(_{j+1,j+1}\) ；因此存在（§ 6 第 1 小节命题 1）E 的正交基 \((e_{j})\) ( \(j=1,\ldots,n\) )，使 \(\Phi(e_{j},e_{j})>0\) （对 \(j=1,\ldots,n\) ）；从而 \(\Phi\) 非退化且正。

注. --- 由标准正交基的存在性可知，两个相同有限维数的向量空间上的两个非退化正埃尔米特型是等价的（§ 1 第 6 小节）。于是设 L 是 A 上一个有限维埃尔米特空间，其度量型非退化且正，V₁ 与 V₂ 是 L 中两个相同维数的线性簇（§ 6 第 6 小节）；由于度量式在 V₁ 与 V₂ 的方向 T₁ 与 T₂ 上的限制、以及在正交子空间 T₁⁰ 与 T₂⁰ 上的限制分别等价，存在 L 的平移空间 T 的酉自同构 \(u\) ，使 \(u(\mathrm{T}_1) = \mathrm{T}_2\) 且 \(u(\mathrm{T}_1^0) = \mathrm{T}_2^0\) ；因而存在 L 的位移 \(v\) 使 \(v(\mathrm{V}_1) = \mathrm{V}_2\) 。设 \((a, b), (a', b')\) 是 L 的点的两个对；要存在 L 的位移 \(v\) 使 \(v(a) = a'\) 且 \(v(b) = b'\) ，必须且只需（用 § 6 第 6 小节的记号）有 \(e(a, b) = e(a', b')\) ；A 的元素 \(\sqrt{e(a, b)}\) （第 VI 章 § 2 第 4 小节）称为埃尔米特空间 L 中 \(a\) 与 \(b\) 的距离。

